\begin{table}[htbp]
  \centering
  \caption{Position of this work relative to nearby scheduling and sensing literatures.}
  \label{tab:related_work_positioning}
  \footnotesize
  \setlength{\tabcolsep}{3.0pt}
  \renewcommand{\arraystretch}{0.96}
  \begin{tabular}{@{}>{\raggedright\arraybackslash}p{0.23\linewidth}>{\raggedright\arraybackslash}p{0.25\linewidth}>{\raggedright\arraybackslash}p{0.22\linewidth}>{\raggedright\arraybackslash}p{0.20\linewidth}@{}}
    \toprule
    Literature line & Typical optimization target & Scheduling structure & What this paper adds \\
    \midrule
    State-estimation and AoI scheduling
      & Covariance, estimation error, information age
      & Sensor selection under communication or energy limits
      & Multi-step forecast loss is the objective. \\
    Active perception and adaptive observation selection
      & Task information under partial observability
      & Observation actions selected from partial observations
      & A mandatory meteorological core channel, limited optional-channel capacity,
        and baselines that use information unavailable during deployment. \\
    DRL-based adaptive sensing and sensor steering
      & Information gain, coverage, data acquisition utility, or monitoring reward
      & Learned adaptive policies, often with soft or task-specific constraints
      & A finite candidate action set enforces the mandatory-core and
        optional-channel-capacity rules, and a feasible-action mask checks the
        steady-state and startup-peak power limits. \\
    Time series forecasting models
      & Forecast accuracy after the observation stream is given
      & Forecasting model for a given observation stream
      & A frozen primary forecaster evaluates every schedule; its parameters are fixed during policy training and evaluation. \\
    \bottomrule
  \end{tabular}
\end{table}
